%% ============================================================
%% Rozdział 2 — Metodologia
%% ============================================================
%% nagłówek rozdziału oraz wariant skrócony do nagłówków/spisu treści
\chapter[Opis metodologii (Imię i Nazwisko autora rozdziału w przypadku pracy zbiorowej)]{Opis metodologii}\label{chapter:2_Methodology} 

W rozdziale tym należy przedstawić szczegółowy opis metod, narzędzi i technik zastosowanych w pracy. Powinien on stanowić rozwinięcie problemów teoretycznych zarysowanych we wstępie (patrz Rozdział \ref{chapter:1_Introduction}) i jasno wskazywać, w jaki sposób autor podszedł do realizacji celu pracy.

W pierwszej kolejności należy opisać zastosowane metody badawcze lub projektowe, ich uzasadnienie oraz kontekst wykorzystania. W przypadku badań empirycznych trzeba wskazać grupę badawczą, sposób doboru próby, narzędzia pomiarowe i procedurę badania. Dla prac projektowych lub inżynierskich należy przedstawić technologie, środowiska programistyczne, biblioteki i frameworki, wraz z krótkim omówieniem przyjętych założeń technicznych.

Następnie należy zaprezentować etapy realizacji projektu w logicznej kolejności – od analizy wymagań i wstępnych koncepcji, poprzez implementację, aż po testowanie i walidację wyników. Wskazane jest uzupełnienie opisu o schematy blokowe, diagramy UML, wzory matematyczne, rysunki lub tabele, które pozwolą na bardziej przejrzyste przedstawienie rezultatów.

\paragraph{Przykładowa struktura metodologii:} 
\begin{enumerate}
    \item Charakterystyka metod i narzędzi – opis wybranych technik badawczych, algorytmów, środowisk programistycznych lub urządzeń.
    \item Opis etapów realizacji – kolejno przedstawione działania, wraz z uzasadnieniem ich zastosowania.
    \begin{enumerate}[label=\alph*)]
        \item Etap 1: analiza wymagań.
        \item Etap 2: projektowanie rozwiązania.
        \begin{itemize}
            \item Etap 3: implementacja i integracja.
            \item Etap 4: testowanie i weryfikacja wyników.
        \end{itemize}
    \end{enumerate}
    \item Podsumowanie – krótkie zestawienie przyjętej metodyki i jej znaczenia dla osiągnięcia celów pracy.
\end{enumerate}

Całość powinna być przygotowana w sposób spójny, tak aby czytelnik mógł jasno prześledzić logikę i konsekwencję działań prowadzących do uzyskania wyników pracy
Ewentualne przypisy można wstawić za pomocą polecenia \verb|\footnote{Treść.}| np. \footnote{Ich oznaczenie to kolejne wartości narastające zgodnie z kolejnością występowania w pracy.}
